\section{Video}
\label{sec:video}

The demonstration video is available at
\url{https://www.youtube.com/watch?v=JrKrHBp4w70}. It is unlisted rather than
private, so the link opens without a YouTube account.

It shows the client connecting to the flashed board and games played through
the GUI: a new game, legal moves each confirmed by the server before the board
is redrawn, a resignation, and a game played out to a win, after which the
cells and the Resign button are both disabled. The board is in frame beside the
screen throughout, with the on-board button and LED visible. A refused move is
not demonstrated in the video; Fig.~\ref{fig:gui-rejected} shows that case
instead, and Section~\ref{sec:valid} asserts it automatically in seven of the
thirty-three tests.

\section{Code}
\label{sec:code}

The source code is at \url{https://github.com/TriTacLe/serial-tictactoe}. The
repository is public, so the link opens without a GitHub account.

The repository holds the board project in \texttt{board/}, the hardware-free
game core and board application in \texttt{server/}, the Tkinter client in
\texttt{client/}, the host simulator in \texttt{sim/} and the pytest harness in
\texttt{tests/}. Setup, build, flash, run and test instructions are in the
top-level \texttt{README.md}, with the board-specific steps in
\texttt{board/README.md}.
